% ---------------------------------------------------------------------
% Preâmbulo: pacotes, geometria, cores, macros e ambientes
% ---------------------------------------------------------------------
\usepackage[brazilian,english]{babel}
\usepackage[utf8]{inputenc}
\usepackage[T1]{fontenc}
\usepackage{amsmath, amssymb, amsfonts}
\usepackage{graphicx}
\usepackage{geometry}
\usepackage{enumitem}
\usepackage{xcolor}
\usepackage{colortbl}
\usepackage{setspace}
\usepackage{array}
\usepackage{longtable}
\usepackage{needspace}
\setlength{\tabcolsep}{5pt}
\usepackage{float}
\usepackage{booktabs}
\usepackage{multirow}
\usepackage{tikz}
\usepackage[hyphens]{url}
\usepackage{hyperref}
\Urlmuskip=0mu plus 1mu\relax

\geometry{top=2.5cm, bottom=2.5cm, left=2.5cm, right=2.5cm}

% ---------- Cores ----------
\definecolor{azulPetCuida}{HTML}{0B5A85}
\definecolor{destaque}{HTML}{DCEAF5}
\definecolor{cinzaclaro}{HTML}{EEF2F6}
\definecolor{verdeP}{HTML}{1B7F3B}
\definecolor{laranjaH}{HTML}{C26A00}
\definecolor{vermelhoC}{HTML}{B3261E}

\hypersetup{
    colorlinks=true,
    allcolors=black,
    pdftitle={Tipo societário e aspectos legais -- PetCuida},
    pdfauthor={PetCuida -- Tópicos III},
}

% ---------- Espaçamento ----------
\setlength{\parindent}{1.25cm}
\setlength{\parskip}{4pt}
\setlength{\emergencystretch}{2em}
\renewcommand{\arraystretch}{1.25}
\setlist{itemsep=2pt, topsep=4pt}

% ---------- Dados do trabalho (edite aqui; reflete na capa e folha de rosto) ----------
\newcommand{\projeto}{PetCuida}
\newcommand{\tituloatividade}{Tipo societário e aspectos legais}
\newcommand{\tipoatividade}{(ATIVIDADE EM GRUPO)}
\newcommand{\anoentrega}{2026}
\newcommand{\orientador}{João Carlos Oliveira Caetano}
\newcommand{\tutor}{Paulo Mattos}
\newcommand{\numeroatividade}{X}
\newcommand{\datareferencia}{15 de outubro de 2026}

% ---------- Ambientes de elementos pré-textuais ----------
\newenvironment{capa}{\thispagestyle{empty}}{\clearpage}
\newenvironment{folhaderosto}{\thispagestyle{empty}}{\clearpage}

% ---------- Colunas e cabeçalhos de tabela ----------
\newcolumntype{L}[1]{>{\raggedright\arraybackslash}p{#1}}
\newcolumntype{C}[1]{>{\centering\arraybackslash}p{#1}}
\newcolumntype{R}[1]{>{\raggedleft\arraybackslash}p{#1}}
\newcommand{\cab}[1]{\textbf{\color{white}#1}}
\newcommand{\cabfundo}{\rowcolor{azulPetCuida}}
\newcommand{\totalfundo}{\rowcolor{destaque}}